\subsection{Rational-tail normalizers in the full completion}

\begin{result}{E.33}{Proved}{All-genus tripod normalizers and their crossing maps}
Start with a stable surface of type $(g,n)$, choose distinct
marked points $p_1,\ldots,p_r$, and replace each by a labelled
pair $(p_i,x_i)$ on an attached three-pointed rational tail.
Let $\sigma$ be the resulting multicurve on the surface of
type $(g,n+r)$, $\Gamma$ its pure mapping class group,
$C=\Stab_\Gamma(\sigma)$, and
$\Delta=\langle\tau_1,\ldots,\tau_r\rangle$ its twist lattice.
Then
\[
 \widehat C\lhook\joinrel\longrightarrow\widehat\Gamma,
 \qquad
 N_{\widehat\Gamma}(\widehat\Delta)
   =C_{\widehat\Gamma}(\widehat\Delta)=\widehat C.             \tag{E.100}
\]
In particular the right action of
$\widehat{C/\Delta}$ on the completed crossing
$\widehat\Gamma/\widehat\Delta$ is faithful and free.

Fix this collection of tails and take the full subcategory
whose objects are the open stratum and the strata associated
to its subsets.  A continuous autoequivalence fixing these
objects is determined, up to natural isomorphism, by its
outer action on $\widehat\Gamma$.  The comparison includes
all its crossing maps, and uses the full completions in every
genus.  For $r=n$ its deepest crossing recovers the fully
boundary-fixed group $\widehat\Gamma_{g,n}^{\partial}$.
\end{result}
\begin{proof}
We first supply the peripheral calculation for every
punctured hyperbolic surface used here.  For its discrete
fundamental group $\pi$, a peripheral cyclic group has its
full topology in $\pi$.  In genus zero its homology class
is primitive.  In positive genus, pass to the cover induced
by reduction modulo two of the compact surface homology.
It is unramified on compactification, has degree at least
four, and contains every peripheral group.  The cover has
at least four punctures, so each of its peripheral homology
classes is primitive.  Finite-index passage proves the
assertion for $\pi$.

Write $I\cong\widehat\Z$ for the corresponding group in
$\widehat\pi$.  We prove
\[
                  N_{\widehat\pi}(I)=I.                       \tag{E.101}
\]
Suppose an element $u$ outside $I$ normalizes it.  Choose
an open subgroup containing $I$ and excluding $u$.
Its corresponding finite cover has at least three punctures
in the genus-zero case.  In positive genus its compactification
has genus at least one by Riemann--Hurwitz; replacing it by
the compact-homology cover just used makes the number of
punctures at least four.  Denote the resulting open subgroup
by $V$.  Both $VI$ and $VuI$ are unramified cusp double
cosets, and they are distinct.  Their inertia subgroups in
$V$, up to conjugacy, are both $I$: a discrete representative
of $Vu$ differs from $u$ by an element of $V$.
On a surface with at least three punctures, the peripheral
homology classes have only their single sum relation.
Two distinct cusps therefore span distinct procyclic lines
in completed homology.  The two identified inertia groups
give the same line, a contradiction.  This proves (E.101),
by the actual covering and homology maps, for every genus.

The full completed Birman sequence needed below is
\[
 1\longrightarrow\widehat\pi_{g,m}\longrightarrow
 \widehat\Gamma_{g,m+1}\longrightarrow
 \widehat\Gamma_{g,m}\longrightarrow1,\qquad m\geq1,           \tag{E.102}
\]
in the stable range.  Here is the completion check.
The discrete exact sequence follows from the evaluation
fibration and the relative identity-component calculation
in C.5.  Its kernel is the point-pushing group of the
punctured surface and is a finite-rank nonabelian free group.
Characteristic finite-index subgroups of that group are
cofinal.  Conjugation by the total mapping class group
therefore gives homomorphisms to their finite automorphism
groups, and hence a continuous action of the full completed
mapping class group on $\widehat\pi_{g,m}$.
On the completed point-pushing subgroup this is the inner
action.  A nonabelian free profinite group has trivial
center: view a free basis as peripheral generators of a
punctured sphere and apply (E.101) to two independent basis
elements.  Their procyclic subgroups have trivial intersection,
as their two exponent coordinates show.
Thus the map of $\widehat\pi_{g,m}$ into the completed
mapping class group is injective.  Completion of the quotient
sequence is right exact, which proves (E.102).

For one rational tail, forgetting $x_1$ sends its stabilizer
onto the base mapping class group and has kernel
$\langle\tau_1\rangle$.  Explicitly the stabilizer is the
group with a framed boundary at $p_1$, and the quotient caps
that boundary by the marked disk.  Representatives fixing
this disk give surjectivity.  Its kernel is the point-push
of $x_1$ around $p_1$, the peripheral generator in (E.102).
For $r$ tails, the same cutting-and-capping maps give
\[
 1\longrightarrow\Delta\longrightarrow C\longrightarrow
                 \Gamma_{g,n}\longrightarrow1.               \tag{E.103}
\]
Every tail is fixed individually because its two labels are
fixed.  Thus $C$ centralizes all the $\tau_i$.

The map on the right of (E.103) is the restriction of the
global map forgetting $x_1,\ldots,x_r$.  In E.29 the full
topology on $C$ was proved using its diagonal map to
$\Gamma\times(C/\Delta)$.  Here the second coordinate
extends to a homomorphism on all of $\Gamma$.
Consequently that product topology on $C$ is already its
topology induced from $\Gamma$.  This proves the injection
in (E.100) and the exact completed version of (E.103).

For $r=1$, take $u\in N_{\widehat\Gamma}(\widehat\Delta)$.
Lift its image under (E.102) to $c\in\widehat C$, using
the completed surjection in (E.103).  Then $c^{-1}u$ lies
in the completed point-pushing group and normalizes the
cusp group $\widehat\Delta$.  Equation (E.101) places it
in that group, so $u\in\widehat C$.

Proceed by induction on $r$.  Inside $\widehat\Delta$,
the kernel of forgetting $x_i$ is exactly
$\langle\tau_i\rangle^{\wedge}$: the other tail twists
retain their full independent lattice on the target.
The global forgetful kernel is normal, so an element
normalizing $\widehat\Delta$ normalizes each of these
individual factors.  Forget $x_r$ and apply the induction
assertion to the other $r-1$ tails.  The image belongs to
the completion of their stabilizer.  The map from $C$
to that stabilizer is onto, with kernel
$\langle\tau_r\rangle$, by capping just its last framed
boundary.  Multiply by a lift in $\widehat C$ to enter
the kernel in (E.102).  The remaining element normalizes
the cusp at $p_r$, so (E.101) puts it in
$\langle\tau_r\rangle^{\wedge}\subset\widehat C$.
This proves the normalizer equality.  Since $C$ centralizes
the twist lattice, its completion does too; the centralizer
equality follows from the two inclusions.

We verify the categorical assertion, including the topology
on its intermediate arrows.  Choose the actual tail curves
simultaneously, so every subset uses its inherited
representative.  The fixed labels force every stabilizer
of a subset to fix each of its tails.  For nested subsets
$\sigma_S\supset\sigma_T$, put
$C_T=\Stab(\sigma_T)$.  The transporter formula then gives
\[
 \Hom(\sigma_S,\sigma_T)=C_T/\Delta_S,
 \qquad
 \Hom(\sigma_S,\varnothing)=\Gamma/\Delta_S.                  \tag{E.104}
\]
Indeed a transporter preserving the indicated target tails
belongs to $C_T$, and every such element is a transporter;
the graph labels exclude any permutation to a different
selected tail.  Each $C_T$ has its full topology induced
from $\Gamma$, by the argument above.

Any finite categorical quotient on the first Hom-set in
(E.104) is a finite transitive $C_T$-set.  Its action kernel
is a finite-index subgroup $U$ of $C_T$.
Choose $N\triangleleft\Gamma$ of finite index with
$N\cap C_T\subset U$.  The ambient finite transporter
quotient (E.82) then refines that Hom quotient, because
\[
       (N\Delta_S)\cap C_T=(N\cap C_T)\Delta_S
                                  \subset U\Delta_S.
\]
This proves cofinality and identifies the two completed
Hom-sets and their comparison as
\[
  \widehat C_T/\widehat\Delta_S
       \ \lhook\joinrel\longrightarrow\
       \widehat\Gamma/\widehat\Delta_S.                      \tag{E.105}
\]
The map is composition with the chosen crossing from
$\sigma_T$ to the open object.  It is injective by (E.100).

Now make the open action of an autoequivalence in the
asserted kernel literally the identity, using its inner
conjugator as the component of a natural isomorphism.
The image of each chosen crossing lies in its homogeneous
space $\widehat\Gamma/\widehat\Delta_S$ and has the same
left stabilizer.  Equation (E.100) says that its displacement
comes from a unique element of
$\widehat{C_S/\Delta_S}$ acting on the right.
Use that element as the component of a further natural
isomorphism.  All chosen crossings are now fixed.
Their right actions are free, since
$\widehat C_S\to\widehat\Gamma$ is injective and the
kernel of $\widehat C_S\to\widehat{C_S/\Delta_S}$ is
$\widehat\Delta_S$.  Naturality consequently fixes every
endpoint automorphism.  It fixes the remaining arrows
by the injectivity of (E.105).  The normalized functor is
the identity, proving the claimed determination.

For $r=n$ the cutting presentation (E.84) gives
$C=\Gamma_{g,n}^{\partial}$ with its literal inverse on
each core and tripod generator.  The topology and all
crossing identifications just proved apply to that map.
\end{proof}
